\section{Reference process and entropy comparison}\label{sec:referencia}

The continuous profile supplies the weights for constructing placements one step at a time. This section
compares the law of the
resulting placements, conditional on success, with the uniform law on the layer. The relative entropy bound
will allow us to transfer
estimates obtained during the construction to the uniform law.

Fix $K$ and $0\le h\le K$, and set \[
 q=n-h,\qquad N=2n-1,\qquad H=N-q=n+h-1 .
\] In this section and the next two, for $h\ge1$ we evaluate the profile at the target parameter $x_*=h/n$,
which is held fixed
throughout the process, and omit the subscript $x_*$. Where this target profile and the full-board profile
appear together, as in
Section~\ref{subsec:trayectoria}, we write the subscript $x_*$ again on the former. We treat layer $h=0$
separately in
Section~\ref{subsec:h0}.

\subsection{Augmented hypergraph and weights}

We add $n-1$ \emph{dummy} vertices to the rows and as many to the columns, so that each of the four families
contains $N$ vertices. The
\emph{real edges} are the $n^2$ squares, each with its four physical lines. The \emph{virtual edges} are all
quadruples containing one
vertex from each family, whether physical or dummy. The color (real or virtual) is part of the edge. A
virtual edge is distinct from a
real edge even if they have the same vertices.

In Figure~\ref{fig:cruces-vacios}, the lines through $(3,2)$ are $R_3$, $C_2$, $D_1$, and $A_5$ because
$3-2=1$ and $3+2=5$. Choosing
the real edge on that square places a queen there. Choosing the virtual edge with the same four lines
deactivates those lines in the
process without placing a queen on the board. We will use this distinction when counting repairs. We call
the edges in a trajectory
\emph{pieces}. These are the edges chosen by the process, real or virtual according to their color.

For $h\ge1$, the vertex weights are $u_R(\bar r)=x_*r_R(\bar r)$ on the physical rows,
$u_C(\bar c)=x_*r_C(\bar c)$ on the physical
columns, $u_D$ and $u_A$ at the induced coordinates of the diagonals and antidiagonals, respectively, and
$1$ on the dummy vertices.
All weights are positive. The weight of an edge, real or virtual, is the product
$\varpi(e)=\prod_{v\in e}u_v$. By
Theorem~\ref{thm:perfil}(a), every real edge has weight between $cx_*^2$ and $Cx_*^2$.

\subsection{The process}

\begin{definicion}[Reference process]\label{def:referencia} Choose in advance, without revealing it, a
uniformly random word with $q$
occurrences of the real color and $H$ of the virtual color. At each step $t=1,\dots,N$, reveal only the next
color, choose an available
edge of that color (one whose four vertices are still active) with probability proportional to its weight,
and deactivate its four
vertices. A virtual piece places no queen. If the normalizer for the requested color, namely the sum of the
weights of its available
edges, is zero, the process enters a failure state. Let $\QQ_h$ be the law of the trajectories and $S$ the
success event (all $N$ steps
are completed). \end{definicion}

On $S$, the $q$ selected real edges are squares with no lines in common. They form a placement in
$\Om_h(n)$, called the \emph{final
placement}. Let $\nu'_h$ be its law under $\QQ_h(\cdot\mid S)$. For real choices the common factor $x_*^2$
can be cancelled without
changing the law; it does not cancel for virtual choices. The reference process need not succeed with
probability close to one.

\subsection{The comparison law}

\begin{definicion}[Comparison law]\label{def:comparacion} Choose $P\in\Om_h(n)$ uniformly at random, then
choose a uniform virtual
completion, that is, a perfect matching by virtual edges of the $H$ free vertices in each family (there are
$(H!)^3$ such matchings).
Finally, choose a uniform ordering of the $N$ pieces. Let $\PP_h$ be the law of the resulting sequence of
edges. \end{definicion}

We represent the uniform order by assigning the pieces independent uniform times in $[0,1]$ and revealing
them in increasing order of
time. (This uses independence. Uniformity of each individual time would not suffice.) These independent
times, together with
Jensen's inequality in Lemma~\ref{lem:normalizadores}(iii), are the device of the entropy bounds of
Radhakrishnan~\cite{Radhakrishnan} and of Linial and Luria~\cite{LinialLuria}. The law $\PP_h$ projects
exactly onto $\mu_h$, and
each trajectory in its support has positive probability under $\QQ_h$ because the next piece of the complete
matching remains available
and has positive weight. In particular, $\PP_h(S)=1$.

\subsection{The entropy identity}

Let $W_R(t)$ and $W_V(t)$ be the total weights of the available edges of the two colors before step $t$, and
let $\col(t)$ denote the
color revealed at that step.

\begin{lema}\label{lem:identidad} For $1\le h\le K$, \[
 \KL{\PP_h}{\QQ_h}=\E_{\PP_h}\sum_{t=1}^N\log W_{\col(t)}(t)-\sum_v\log u_v-\log N_h-\log q!
 -4\log H!,
\] where $\sum_v$ ranges over the $4N$ vertices. \end{lema}

\begin{proof} Let $\omega$ be a trajectory in the support of $\PP_h$. Then
$\PP_h(\omega)=[N_h(H!)^3N!]^{-1}$. Under $\QQ_h$, the color
word has probability $\binom Nq^{-1}$ and each choice has probability $\varpi(e_t)/W_{\col(t)}(t)$. Since
the complete matching covers
every vertex exactly once, the product of the chosen weights is $\prod_vu_v$, independent of the matching
and its order. Hence \[
 \log\frac{\PP_h(\omega)}{\QQ_h(\omega)}=\sum_t\log W_{\col(t)}(t)-\sum_v\log u_v
 -\log\bigl[N_h(H!)^3N!\bigr]+\log\binom Nq ,
\] and $-\log[(H!)^3N!]+\log\binom Nq=-\log q!-4\log H!$. Taking the expectation under $\PP_h$ gives the
result. \end{proof}
